\documentclass[pdftex,a4paper]{article}
\usepackage[lmargin=1.25in,rmargin=1.25in,tmargin=1in,bmargin=1in]{geometry}
\usepackage[T1]{fontenc}
\usepackage{lmodern}
\usepackage[utf8]{inputenc}
\usepackage[czech]{babel}
\usepackage{url}
\usepackage{graphicx}
\usepackage[usenames,dvipsnames]{xcolor}
\usepackage{listings}
\usepackage[unicode,pdfdisplaydoctitle,pdftex,colorlinks,allcolors=black]{hyperref}
\hypersetup {
    pdftitle={Vizualizace grafu matematické funkce},
    pdfauthor={Vojtěch Kinkor},
}
\title{Vizualizace grafu matematické funkce} 
\author{Vojtěch Kinkor}     
\date{\today}   

\lstset{ %
   backgroundcolor=\color{white},   
   basicstyle=\footnotesize\ttfamily\fontfamily{pcr}\selectfont,      
   breakatwhitespace=true,          		
   breaklines=true,                 		
   commentstyle=\itshape\color{Gray},    	
   keywordstyle=\bfseries\color{blue},      
   stringstyle=\color{ForestGreen},     	
   escapeinside={\%*}{*)},          		
   frame=none,                    			
   keepspaces=true,                    
   language=C,                 				
   numbers=none,                    		
}


\begin{document}

\begin{titlepage}
\begin{center}


~\\[2.5cm]
\includegraphics[width=0.7\textwidth]{res/logo}
\\[0.5cm]


{ \huge \bf Semestrální práce z předmětu\\ KIV/SU \\[0.6cm] }
{ \LARGE \sc Shlukovací algoritmy }

\vfill

% Author and supervisor
\begin{minipage}{0.49\textwidth}
\begin{flushleft} 
\emph{Datum:}\\
\today
\end{flushleft}
\end{minipage}
\begin{minipage}{0.49\textwidth}
\begin{flushright} 
\emph{Vypracovali:} \\
Matěj \textsc{Kareš}\\
Vojtěch \textsc{Kinkor}
\end{flushright}
\end{minipage}



\end{center}
\end{titlepage}

\newpage
\tableofcontents




\newpage
\section{Zadání}

Naprogramujte v~ANSI C přenositelnou
\footnote{Je třeba, aby bylo možné Váš program přeložit a~spustit na PC s~operačním prostředím Win32/64 (tj. operační systémy Microsoft Windows NT/2000/XP/Vista/7/8) a~s~běžnými distribucemi Linuxu (např. Ubuntu, OpenSUSE, Gentoo, Debian, atd.). Server, na který budete Vaši práci nahrávat a~který ji otestuje, má nainstalovaný operační systém Debian GNU/Linux 7.1 Wheezy s~jádrem verze 2.6.26-2 a~s~překladačem gcc 4.4.5.} 
\textbf{konzolovou aplikaci}, která jako vstup načte z~parametru na příkazové řádce matematickou funkci ve tvaru y = f(x), provede její analýzu a~vytvoří soubor ve formátu \textsf{PostScript} s~grafem této funkce na zvoleném definičním oboru.
\\\\
Program se bude spouštět příkazem \texttt{graph.exe $\langle$func$\rangle$ $\langle$out-file$\rangle$ [$\langle$limits$\rangle$]}. Symbol \texttt{$\langle$func$\rangle$}
zastupuje zápis dvojdimenzionální matematické funkce (funkce ve více dimenzích program řešit nebude, nalezne-li během analýzy zápisu funkce více nezávislých proměnných než jednu, vypíše srozumitelné chybové hlášení a~skončí). Symbol \texttt{$\langle$out-file$\rangle$} zastupuje jméno výstupního postscriptového souboru. Takže Váš program může být během testování spuštěn například takto:
\\[2mm]
\texttt{\ldots\char`\\ >graph.exe\ "sin(2*x)\^{}3"\ mygraph.ps}
\\[2mm]
Výsledkem práce programu bude soubor ve formátu \textsf{PostScript}, který bude zobrazovat graf zadané matematické funkce – ve výše uvedeném případě $y =\ $sin$(2x)^3$ – v~kartézské souřadnicové soustavě $(O; x, y)$ s~vyznačenými souřadnými osami a~(aspoň) význačnými hodnotami definičního oboru a~oboru hodnot funkce.
\\\\
Pokud nebudou na příkazové řádce uvedeny alespoň dva argumenty, vypište chybové hlášení a~stručný návod k~použití programu v~angličtině podle běžných zvyklostí (viz např. ukázková semestrální práce na webu předmětu Programování v~jazyce C). \textbf{Vstupem programu jsou pouze argumenty na příkazové řádce – interakce s~uživatelem pomocí klávesnice či myši v~průběhu práce programu se neočekává.}
\\\\
Hotovou práci odevzdejte v~jediném archivu typu ZIP prostřednictvím automatického odevzdávacího a~validačního systému. Archiv nechť obsahuje všechny zdrojové soubory potřebné k~přeložení programu, \textbf{makefile} pro Windows i Linux (pro překlad v~Linuxu připravte soubor pojmenovaný \texttt{makefile} a~pro Windows \texttt{makefile.win}) a~dokumentaci ve formátu PDF vytvořenou v~typografickém systému \TeX, resp. \LaTeX. Bude-li některá z~částí chybět, kontrolní skript Vaši práci odmítne.
\\\\\\
Celé zadání: \url{http://www.kiv.zcu.cz/studies/predmety/pc/doc/work/sw2013-02.pdf}






\newpage
\section{Shluková analýza}

Prvotní úvaha se týká celkového postupu pro vyřešení úlohy. Na počátku máme pouze matematickou funkci s~jedním parametrem (resp. \uv{neznámou}). Aby bylo možné vytvořit graf, je potřeba tuto funkci co nejefektivněji vyhodnocovat pro různé parametry. V úloze se tedy jedná ve své podstatě o převedení funkce do jednoduše vyhodnotitelné podoby, vytvoření \uv{kalkulátoru} a~následně opakované vyhodnocování pro vykreslené velkého množství bodů do grafu.
\\\\
Úlohu lze tedy snadno rozdělit do několika samostatných částí:
\begin{itemize}
\item Analýza zadané matematické funkce
\item Zpracování a~vyhodnocení matematického výrazu
\item Vytvoření souboru ve formátu \textsf{PostScript} s~grafem 
\end{itemize}


\subsection{K-means}\label{subsect:analyza}

Matematická funkce se do programu předává pomocí parametru na příkazové řádce a~program s~ní nadále pracuje jako s~řetězcem znaků. Můžeme tedy číst vstup znak po znaku. Ve funkci se mohou obecně objevit tyto symboly (v závorce vždy znaky, ze kterých se symbol může skládat):
\begin{itemize}
\item Konstanty (\texttt{0}--\texttt{9 . E e + -})
\item Proměnné (\texttt{A}--\texttt{Z} \texttt{a}--\texttt{z} -- jeden znak)
\item Aritmetické operátory (\texttt{+ - * / \^{}})
\item Funkce (\texttt{A}--\texttt{Z} \texttt{a}--\texttt{z} -- dva a~více znaků)
\item Závorky (\texttt{( )})
\end{itemize}~

Při zpracování je potřeba jednotlivé symboly od sebe oddělit a~také odhalit různé chyby v~zadaném výrazu. Postupným přidáváním částí pro jednotlivé symboly jsem dospěl k~následujícímu postupu, který se mi jeví jako poměrně nenáročný, ale přesto robustní:

\begin{enumerate}
\item Načtu jeden znak ze vstupního řetězce.
\item Žádný ze symbolů nepožaduje výhradně malá nebo velká písmena, proto pokud je znak velké písmeno, převedu jej na malé.

\item Načtený znak odpovídá sadě znaků pro konstanty -- znaky \texttt{0}--\texttt{9}; \texttt{e}, pokud předcházelo číslo; \texttt{+ -}, pokud předcházel znak \texttt{e}:
\begin{itemize}
\item Pokud není uložený začátek konstanty v~řetězci (tj. pozice prvního znaku), tak uložím aktuální pozici; jinak se vrátím zpět na krok 1. a~pokračuji ve čtení dalšího znaku.
\end{itemize}
\item Pokud podmínka z~bodu 3. neplatí a~je uložený začátek konstanty v~řetězci: 
\begin{itemize}
\item Pak jsem narazil na konec konstanty. Lze ji vykopírovat ze vstupního řetězce a~poté převést na číslo. To následně uložím a~vynuluji údaj pro začátek konstanty.
\end{itemize}

\item Načtený znak odpovídá sadě znaků pro funkce nebo proměnné -- znaky \texttt{a}--\texttt{z}:
\begin{itemize}
\item Pokud není uložený začátek funkce v~řetězci, tak si uložím aktuální pozici; jinak se vrátím zpět na krok 1. a~pokračuji ve čtení dalšího znaku.
\end{itemize}
\item Pokud podmínka z~bodu 5. neplatí a~je uložený začátek funkce v~řetězci:
\begin{itemize}
\item Zpracovaný je pouze jeden znak -- jedná se o proměnnou. Lze ověřit platnost proměnné -- v~tomto případě pouze zjistím, zda se jedná o znak \texttt{x}. Proměnnou uložím.
\item Zpracovaných je více znaků -- poté se jedná o funkci. Ze vstupního řetězce vykopíruji samostatnou funkci a~ověřím, zda taková funkce existuje. Funkci uložím.
\item Nakonec vynuluji údaj pro začátek funkce/proměnné.
\end{itemize}

\item Načtený znak odpovídá sadě znaků pro operátory nebo závorky -- znaky \texttt{+ - * / \^{} ( )}
\begin{itemize}
\item Lze provést kontroly dle různých pravidel (dva operátory za sebou, chybějící závorky, apod.). Podle vstupu uložím nebo jinak zpracuji.
\end{itemize}

\item Pokud načtený znak neodpovídá žádné skupině a~zároveň se nejedná o mezeru nebo ukončovací znak řetězce, pak se jedná o neplatný znak - zpracování se přeruší.

\item Nakonec se proces vrátí zpět na 1. krok a~pokračuje ve čtení dalšího znaku, dokud nenarazí na ukončovací znak.

\end{enumerate}


\subsection{Zpracování a~vyhodnocení matematického výrazu}\label{subsect:zprac-vyhodn}
Jakmile umím zadanou matematickou funkci analyzovat a~tedy i oddělit jednotlivé symboly, je potřeba ji převést do vyhodnotitelné podoby. Tato podoba má velký dopad na výslednou efektivitu během vyhodnocování a~tedy i rychlost celé aplikace. 

Rozhodoval jsem se mezi dvěma běžně používanými metodami -- přímou tvorbou binárního stromu a~tvorbou jednorozměrného seznamu s~výrazem v~postfixové notaci (také známá jako \textit{RPN notace}). 

Během realizace jsem tvorbu binárního stromu zavrhl ze dvou důvodů -- samotné vytvoření vyžaduje poměrně složitou analýzu. Alternativně je možné výraz nejprve převést do postfixové notace a~až poté strom vytvářet. Dalším důvodem je ale větší výpočetní náročnost, kvůli možné rozvětvenosti stromu. Procházení jednorozměrného \uv{plochého} seznamu v~druhé metodě bude vždy méně náročné.

Druhá, nakonec zvolená, metoda spočívá v~implementaci algoritmu Shunting-Yard\footnote{algoritmus vyvinutý E. Dijkstrou. Popis viz: \url{http://en.wikipedia.org/wiki/Shunting-yard_algorithm}} do analýzy zadané funkce, čímž získáme výraz v~postfixové notaci uložený právě v~jednorozměrném seznamu. K vyhodnocení se využívá zásobník a~samotný algoritmus je značně jednoduchý:  

\begin{enumerate}
\item Načtu jeden znak.
\item Pokud je znak konstanta -- vložím ji do zásobníku.
\item Pokud je znak funkce nebo operátor potřebující \textit{n} parametrů:
\begin{itemize}
\item Zjistím, zda mám v~zásobníku alespoň \textit{n} prvků -- pokud ne, je zadaný výraz neplatný.
\item Jinak ze zásobníku vyberu právě \textit{n} prvků, následně je pomocí příslušné funkce vyhodnotím. Výsledek vložím zpět do zásobníku.
\end{itemize}
\item Vrátím se na krok 1., pokud jsou ve výrazu další znaky.
\item Jestliže zbylo v~zásobníku více než 1 prvek, je zadaný výraz neplatný.
\item V opačném případně je zbylý prvek v~zásobníku výsledkem celého výrazu.
\end{enumerate}


\subsection{Vytvoření souboru ve formátu \textsf{PostScript} s~grafem}\label{subsect:vytv-ps}
V této části není moc problémů, které by bylo potřeba řešit na teoretické úrovni. Základem je ruční vytvoření ukázkového grafu a~následné \uv{propojení} s~výpočetní částí programu. Jako předloha mi posloužil obrázek grafu, který se nachází v~zadání úlohy. 

Za nejsložitější považuji správné vytvoření popisků os. Na začátku je potřeba zjistit maximální rozsahy pro obě osy. Ty se počítají jako absolutní hodnota z~rozdílu zaokrouhlených hodnot výstupního rozsahu grafu (min. hodnota se zaokrouhluje dolů, max. naopak nahoru) -- tj. pokud má mít osa X hodnoty od 3.14 do 6.28, pak bude její rozsah 4. Poté lze vypočíst přesný počet \uv{značek} na ose a~jejich odstup (\uv{rozměrový} v~grafu i hodnotový pro popisky) a~podle těchto hodnot je mohu zakreslit. Hodnotový rozestup značek by měl být co nejčitelnější (na jedno desetinné místo do rozsahu 5, jinak na celá čísla).

Pro zakreslování hodnot funkce je nutné zvolit vhodnou \uv{jemnost} grafu, což odpovídá počtu hodnot, které se vypočtou a~následně do grafu zakreslí pro zadaný rozsah. Graf také nemusí být spojitý. Klasické vyšetření spojitosti není triviální záležitostí. Při dostatečně \uv{jemném} grafu je však možné jej nahradit velmi jednoduchým ověřování blízkého okolí aktuálního bodu, kdy jen hledáme příliš velký rozdíl mezi hodnotami v~předchozím, aktuálním a~následujícím bodě.


\newpage
\section{Popis implementace}

\subsection{Struktura a~řídící modul \texttt{main.c}}
Program se skládá z~řídícího modulu \texttt{main.c} a~dalších 6 modulů.

Řídící modul zajišťuje především zpracování vstupních argumentů z~příkazové řádky, měření času běhu programu, případně spuštění nestandardních režimů (viz část \ref{subsect:nestd-rezimy} na straně \pageref{subsect:nestd-rezimy}). Nepovinný parametr s~rozsahem výstupního grafu se rozpoznává pomocí funkce \texttt{sscanf}, která umožňuje načítat dle vstupní masky z~řetězců. Další nepovinné parametry se porovnávají funkcí \texttt{strcmp}.

\subsection{Modul \texttt{list.c}/\texttt{list.h}}
Tento modul zastupuje dynamický datový seznam pro uchovávání různých datových objektů, v~tomto případě matematických symbolů. Základní části je opakující se datová struktura:

\begin{lstlisting}
typedef struct List {
    struct List *prev;
    struct List *next;  
    struct MathObject *value;
    int index;
} List;
\end{lstlisting}
%
Tato struktura uchovává odkaz na následující a~předcházející prvek v~seznamu, dále vlastní data -- odkaz na strukturu \texttt{MathObject} -- a~index prvku v~seznamu, který v~praxi slouží pouze ke zjištění celkového počtu prvků.
\\\\
Modul obsahuje tyto funkce:

\begin{itemize}
\item \texttt{List *List\_create()} \\ Vytvoří prázdnou strukturu (seznam, případně prvek seznamu) a~vrátí na ni odkaz, nebo \texttt{NULL}, pokud ji nelze vytvořit.

\item \texttt{void List\_destroy(List *L)} \\ Odstraní stukturu předanou v~parametru.

\item \texttt{int List\_isEmpty(List *L)} \\ Zjistí, zda je zadaný seznam prázdný.

\item \texttt{List *List\_add(List *L, MathObject *value)} \\ Přidá na konec seznamu \texttt{L} prvek s~obsahem \texttt{value}. Vrací odkaz na přidaný prvek, nebo \texttt{NULL}, pokud ji nelze přidat.

\item \texttt{MathObject *List\_removeFirst(List *L)} \\ Vrátí obsah (data) z~prvního prvku seznamu a~poté jej ze seznamu odstraní. Vrací \texttt{NULL}, pokud je seznam prázdný.

\item \texttt{MathObject *List\_removeLast(List *L)} \\ Vrátí obsah z~posledního prvku seznamu a~poté jej ze seznamu odstraní. Vrací \texttt{NULL}, pokud je seznam prázdný.

\item \texttt{MathObject *List\_getFirst(List *L)} \\ Vrátí obsah z~prvního prvku seznamu. Vrací \texttt{NULL}, pokud je seznam prázdný.

\item \texttt{MathObject *List\_getLast(List *L)} \\ Vrátí obsah z~posledního prvku seznamu. Vrací \texttt{NULL}, pokud je seznam prázdný.

\item \texttt{int List\_getNext(List *L, List **N, MathObject **out)} \\ Slouží k~iteraci seznamem. \texttt{L} je odkaz na celý seznam, \texttt{N} odkaz na aktuální prvek (při prvním použití je třeba použít znovu \texttt{L}), \texttt{out} slouží pro předání obsahu aktuálního prvku. Vrací 1, pokud se povedlo přečíst další prvek, jinak 0 (např. na konci seznamu).

\item \texttt{int List\_size(List *L)} \\ Zjistí počet prvků v~seznamu (index posledního prvku).

\item \texttt{int Lists\_checkAlloc(List *lists[], int cnt)} \\ Pomocná funkce -- ověří alokaci struktur předaných v~poli \texttt{lists}. Pokud se některou strukturu nepodařilo naalokovat (odkazuje na \texttt{NULL}), vypíše chybovou hlášku a~ostatní struktury v~poli dealokuje. Vrací 1 v~případě chyby, jinak 0.

\item \texttt{int Lists\_destroy(List *lists[], int cnt)} \\ Pomocná funkce -- smaže zadaný počet struktur předaných v~poli \texttt{lists}.

\end{itemize}




\subsection{Modul \texttt{mathobject.c}/\texttt{mathobject.h}}
Modul obsahující strukturu sloužící k~uchování matematické symbolu (\uv{objektu}) -- odkaz na operaci/funkci, konstantu, nebo neznámou (pozn.: jedná se jen o pozici neznámé v~poli všech používaných neznámých -- dle zadání však toto pole bude vždy obsahovat jen jednu neznámou a~to \texttt{x}).

\begin{lstlisting}
typedef struct MathObject {
    char type;
    union {
        MathOperation *operation;
        double operand;
        int variable;
    } values;
} MathObject;
\end{lstlisting}
%
V proměnné \texttt{type} se uchovává typ \uv{objektu}, který je uložený ve \texttt{values}. 
\\\\
Modul obsahuje tyto funkce:

\begin{itemize}
\item \texttt{MathObject *MathObject\_create\_operation(MathOperationType operation)} \\ 
	  \texttt{MathObject *MathObject\_create\_operand(double operand)} \\
	  \texttt{MathObject *MathObject\_create\_variable(int variable)} \\
	  Vytvoří strukturu s~příslušným obsahem.
	  
\item \texttt{int MathObject\_isOperation(MathObject *M)} \\ 
	  \texttt{int MathObject\_isOperand(MathObject *M)} \\
	  \texttt{int MathObject\_isVariable(MathObject *M)} \\
	  Zjistí, zda je zadaná struktura \texttt{M} příslušného typu.
	  	  
\item \texttt{MathOperation *MathObject\_getOperation(MathObject *M)} \\ 
	  \texttt{double MathObject\_getOperand(MathObject *M)} \\
	  \texttt{int MathObject\_getVariable(MathObject *M)} \\
	  Vrátí obsah ze zadané struktury.
	  	  	  
\item \texttt{MathOperation *MathObject\_removeOperation(MathObject *M)} \\ 
	  \texttt{double MathObject\_removeOperand(MathObject *M)} \\
	  \texttt{int MathObject\_removeVariable(MathObject *M)} \\
	  Vrátí obsah ze zadané struktury a~poté ji odstraní.
	  
\item \texttt{void MathObject\_destroy(MathObject *M)} \\ Odstraní zadanou strukturu.

\end{itemize}



\newpage
\subsection{Modul \texttt{mathoperation.c}/\texttt{mathoperation.h}}
Modul uchovávající funkce, které slouží ke zpracování aritmetických operátorů a~matematických funkcí ve výrazech.

\begin{lstlisting}
#define M_PI 3.14159265358979323846264338327
#define MATHOPERATION_MAXARGS 2

typedef enum { MOFS_LEFT, MOFS_RIGHT, MOFS_NONE, MOFS_CONST } MathOperationFnSide; 
typedef enum { MOR_NAN, MOR_OK, MOR_ERR } MathOperationResult; 
typedef enum { MOT_LB, MOT_RB, MOT_SP, MOT_ADD, MOT_SUB, MOT_MUL, MOT_DIV, MOT_POW, MOT_MOD, MOT_SIN, MOT_COS, MOT_TAN, MOT_ASIN, MOT_ACOS, MOT_ATAN, MOT_SINH, MOT_COSH, MOT_TANH, MOT_LN, MOT_LOG, MOT_EXP, MOT_ABS, MOT_CEIL, MOT_FLOOR, MOT_PI, MOT__LAST } MathOperationType; 

typedef struct MathOperation {
    MathOperationType type;
    unsigned short priority;
    const char *label;
    MathOperationFnSide side;
    short fnArguments;

    MathOperationResult (*compute) (double[], double *);
} MathOperation;

MathOperation MathOperations[MOT__LAST];
\end{lstlisting}
%
Definovaná konstanta \texttt{MATHOPERATION\_MAXARGS} uchovává maximální počet parametrů, které může jakákoliv matematická funkce v~programu zpracovávat. Výčtový typ \texttt{MathOperationFnSide} slouží pro označení asociativity operátorů a~funkcí. \texttt{MathOperationResult} předává informaci o výsledku výpočtu -- \textit{neznámý výsledek}, \textit{OK} a~\textit{chyba během zpracování}. \texttt{MathOperationType} slouží jako označení pro veškeré implementované matematické funkce a~operátory. Struktura \texttt{MathOperation} poté slouží k~uchování informací o dostupném operátoru/funkci i s~odkazem na programovou funkci zajišťující samotné zpracování během výpočtu. \\Globální pole \texttt{MathOperations} seskupuje všechny známé aritmetické operátory a~funkce.
\\\\
Modul obsahuje tyto funkce:

\begin{itemize}
\item \texttt{static MathOperationResult mop\_add (double a[], double *res)} \\
      \ldots \\
	  Skupina (celkem 22) programových funkcí zajišťující samotné výpočty. Předává se pole s~argumenty pro funkci a~odkaz na proměnnou, do které bude výsledek uložen. Funkce vrací výsledek operace.
	  
\item \texttt{MathOperationResult ResultCheck(double *res)} \\ 
	  Pomocná funkce -- jednoduché ověření výsledku operace.
	  	  
\item \texttt{MathOperationType MathOperation\_find(char *op, int len)} \\ 
	  Pokusí se nalézt příslušný operátor/funkci podle názvu. Předává se odkaz na začátek názvu matematické funkce v~řetězci a~jeho délka (počet znaků). Vrací označení příslušné funkce, případně \texttt{MOT\_\_LAST}, pokud nebyla nalezena.\\
	  Pro vyhledávání používá funkci \texttt{strnicmp}, která je ale dostupná pouze pro OS Linux. Při kompilaci pro OS Microsoft Windows se pomocí preprocesoru nadefinuje tato funkce tak, aby byla volána ekvivalentní funkce \texttt{strncasecmp}.
	  	  	  
\end{itemize}



\subsection{Modul \texttt{parser.c}/\texttt{parser.h}}
Modul slouží ke zpracování zadané matematické funkce z~příkazové řádky. Obsahuje jedinou funkci:

\begin{itemize}
\item \texttt{List *ParseMathExpressionToRPN(char *input, char *variables)} \\
	  Pokusí se zpracovat zadanou matematickou funkci v~řetězci \texttt{input}. Parametr \texttt{variables} obsahuje pole známých neznámých. Funkce vrací seznam symbolů v~postfixové notaci. Využívá algoritmus Shunting-Yard. Podrobný popis funkčnosti viz část \ref{subsect:analyza} na straně \pageref{subsect:analyza}. 
 
\end{itemize}


\newpage
\subsection{Modul \texttt{rpncalc.c}/\texttt{rpncalc.h}}
Modul slouží k~vyhodnocení matematického výrazu v~postfixové notaci. Opět obsahuje jedinou funkci:

\begin{itemize}
\item \texttt{MathOperationResult RPNCalc(List *input, double *result, double values[])} \\
	  Vyhodnocuje výraz uložený v~seznamu \texttt{input}, výsledek uloží do proměnné, na kterou odkazuje \texttt{result}. Hodnoty neznámých (v našem případě pouze jedna) se předávají v~poli \texttt{values}. Funkce vrací výsledek výpočtu. Popis algoritmu viz část \ref{subsect:zprac-vyhodn} na straně \pageref{subsect:zprac-vyhodn}. Pro zlepšení výkonnosti funkce se jako zásobník využívá obyčejné pole a~proměnná uchovávající počet prvků v~daném poli. Velikost se stanoví z~počtu prvků v~seznamu \texttt{input} (jelikož se nyní zapisují pouze konstanty, bude toto pole vždy o málo větší, než je potřeba).
 
\end{itemize}


\subsection{Modul \texttt{grapher.c}/\texttt{grapher.h}}
Modul zajišťující samotné vytvoření souboru ve formátu \textsf{PostScript} s~grafem zadané matematické funkce. Obsahuje tři globální celočíselné proměnné -- \texttt{GRAPHER\_WIDTH}, \texttt{GRAPHER\_HEIGHT} a~\texttt{GRAPHER\_PREC} -- sloužící k~nastavení rozměrů výstupního grafu, respektive \uv{jemnosti} vykreslování. Dále modul obsahuje tyto funkce:

\begin{itemize}
\item \texttt{void Grapher\_render(char *input, char *filename, double xmin, double xmax,\\double ymin, double ymax)} \\
	  Funkce, která sama zajistí celé zpracování -- jako parametry se předává řetězec s~matematickou funkcí, název výstupního souboru a~rozsah grafu. Důležité části funkce jsou obecně popsány v~části \ref{subsect:vytv-ps} na straně \pageref{subsect:vytv-ps}. Nejprve se nechá zadaný výraz zpracovat do postfixové notace. V případě úspěchu se začne vykreslovat graf.
	  
\item \texttt{double roundd(double val, double prec)} \\
	  Pomocná funkce -- slouží k~zaokrouhlení čísla \texttt{val} na určitý počet desetinných míst. 
	  	  
\item \texttt{void printList(List *L, char *variables)} \\
	  Pomocná funkce používaná pouze při vývoji -- vytiskne obsah seznamu se zpracovaným matematickým výrazem. Jako druhý parametr se předává pole neznámých.
	  
\end{itemize}


\newpage
\section{Uživatelská příručka}

\subsection{Překlad programu}
Program je možné přeložit nástrojem \texttt{make}\footnote{Ve většině distribucí OS Linux bývá nástroj \texttt{make} předinstalován. Pro OS Microsoft Windows je nutné\\ nainstalovat např. prostředí Cygwin, viz: \url{http://www.cygwin.com/}}, stačí ve složce se zdrojovými soubory zavolat příkaz:
\\[2mm]
\texttt{\ldots\char`\\ > make}
\\[2mm]
Pokud vše proběhne správně, bude ve stejné složce i spouštěcí soubor s~názvem \texttt{graph.exe}.

Pro překlad na OS Microsoft Windows je též možné využít překladové prostředí Visual Studio. Je nutné spustit konzoli \textit{Visual Studio Command Prompt} a~ve složce se zdrojovými soubory zavolat příkaz:
\\[2mm]
\texttt{\ldots\char`\\ > nmake -f Makefile.win}

\subsection{Spuštění programu}
Ovládání programu je přesně popsané v~kompletním zadání úlohy\footnote{\url{http://www.kiv.zcu.cz/studies/predmety/pc/doc/work/sw2013-02.pdf}}. Program v~základní podobě se spouští z~příkazové řádky příkazem\footnote{Pro unixové systémy (např. OS Linux) je nutné na začátek příkazu připsat \texttt{./}}
\\[2mm]
\texttt{graph.exe $\langle$func$\rangle$ $\langle$out-file$\rangle$ [$\langle$limits$\rangle$]}
\\[2mm]
Kde symbol \texttt{$\langle$func$\rangle$} zastupuje matematickou funkci, symbol \texttt{$\langle$out-file$\rangle$} jméno výstupního postscriptového souboru a~nepovinný symbol \texttt{$\langle$limits$\rangle$} definuje rozsah kresleného grafu. Pokud parametry nesouhlasí, vypíše se stručná uživatelská nápověda programu na obrazovku, viz obrázek \ref{fig:obr1}.
\\

\begin{figure}[h]
\includegraphics[width=1\textwidth]{res/2013-12-31_154800.png}
\caption{Výpis uživatelské nápovědy}
\label{fig:obr1}
\end{figure}

\newpage

O úspěšném průběhu je uživatel informován, rovněž tak, pokud se vyskytne chyba -- nejčastěji se jedná o chybu v~zadané matematické funkci. Viz obrázek \ref{fig:obr2}. 

\begin{figure}[h]
\includegraphics[width=1\textwidth]{res/2013-12-31_154938.png}
\caption{Výpis chybové hlášky}
\label{fig:obr2}
\end{figure}


\subsubsection*{Matematická funkce}
Matematická funkce slouží k~popisu 2D grafu, očekává se v~obvyklém zápisu. Může obsahovat konstanty v~libovolném zápisu, proměnnou \texttt{x}, aritmetické operátory, funkce a~závorky. Navíc je přidána konstanta \texttt{pi}.
\\\\
Implementovány jsou následující funkce (\# zastupuje parametr funkce -- matematický výraz):
\begin{itemize}
\item \texttt{sin(\#) cos(\#) tan(\#)} -- goniometrické funkce
\item \texttt{asin(\#) acos(\#) atan(\#)} -- cyklometrické funkce 
\item \texttt{sinh(\#) cosh(\#) tanh(\#)} -- hyperbolometrické funkce 
\item \texttt{ln(\#) log(\#)} -- přirozený a~dekadický logaritmus
\item \texttt{exp(\#)} -- funkce $e^x$
\item \texttt{abs(\#)} -- absolutní hodnota
\item \texttt{ceil(\#) floor(\#)} -- zaokrouhlení nahoru a~dolů 
\item \texttt{mod(\#,\#)} -- zbytek po dělení (\uv{modulo}) 
\end{itemize}
Program umožňuje i zadání dalších dobrovolných parametrů:
\begin{itemize}
\item \texttt{-t} -- Zobrazí dobu běhu programu
\item \texttt{-p \#} -- Nastaví \uv{jemnost} vykreslovaného grafu (implicitně 10, min. 1)
\item \texttt{-w \#}  -- Šířka grafu v~souboru \textsf{PostScript} (implicitně 595)
\item \texttt{-h \#}  -- Výška grafu v~souboru \textsf{PostScript} (implicitně 500)
\end{itemize}

\subsubsection*{Nestandardní režimy běhu programu}\label{subsect:nestd-rezimy}
Program lze spustit pomocí následujících příkazů v~dalších dvou nestandardních režimech, které nebyly v~zadání úlohy požadovány:

\begin{description}
\item[\texttt{graph.exe -calc}] -- Spustí režim interaktivního kalkulátoru. Zadávat lze výrazy ve stejné podobě, jako se používají pro vykreslení grafu. Do proměnné \texttt{x} se vždy ukládá výsledek poslední operace. Napsáním příkazu \texttt{exit} se program ukončí.
\item[\texttt{graph.exe -test}] -- Spustí výkonnostní test s~předdefinovanou funkcí a~počtem opakování. Slouží pro porovnání výkonnosti aplikace při použití různých kompilátorů.
\end{description}







\newpage
\section{Závěr}

Práci jsem vyvíjel na počítači s~procesorem Intel Core i5 M480 s~6GB RAM s~OS Microsoft Windows 7. Používal jsem především překladač, který je součástí vývojového prostředí Microsoft Visual Studio 2010, a~zároveň testoval v~prostředí Cygwin s~překladačem \texttt{gcc} ve verzi 4.8.2 a~na Linuxovém serveru běžícím na platformě ARM s~překladačem \texttt{gcc} 3.4.6.

Doba běhu programu se při implicitním nastavení pohybuje v~řádu desítek ms (20--80ms). Velikost výstupního souboru se běžně pohybuje v~rozmezí 80--150kB. Paměťová náročnost závisí na složitosti zadané matematické funkce -- obecně lze říci, že každý symbol ve funkci odpovídá 50--150B alokované paměti. Při běhu programu v~testovacím režimu, ve kterém se počítá velmi složitá funkce, nepřesáhla alokovaná paměť 2MB a~tato hodnota je po celou dobu programu téměř konstantní. Měření bylo prováděno pomocí nástroje Valgrind.

Během programování jsem si oživil znalosti a~zvyklosti jazyka ANSI C, který považuji za značně programátorsky neefektivní, na druhou stranu jeho striktnost je dobrým vodítkem k~pochopení fungování programů, především pak správy paměti. Novou zkušeností bylo pro mne seznámení s~jazyky \textsf{PostScript} a~sázecím programem \TeX. Oba dva považuji za velmi zajímavé nástroje, především práce s~PostScriptem mě bavila. Psaní dokumentace v~\TeX u není nijak přehnaně složité, některé části jako například tvorba obsahu jsou příjemně zautomatizované a~autor se o ně nemusí nadále starat. 

Program splňuje veškeré body zadání a~s~výsledkem jsem poměrně spokojený. Možná vylepšení vidím v~části kreslení grafu, především pak hledání nespojitostí funkcí a~minimalizaci výstupu. Všechny části programu jsou snadno znovupoužitelné a~rozšiřitelné.





\end{document}












